% ═══════════════════════════════════════════════════════════════════════════
% LERNBLATT: Elektromagnetische Induktion (Kl. 9, DS 09b)
% Begleitend zum digitalen Lernpfad LP-09b
% ═══════════════════════════════════════════════════════════════════════════

\documentclass[a4paper,11pt]{article}

\usepackage[utf8]{inputenc}
\usepackage[T1]{fontenc}
\usepackage[ngerman]{babel}
\usepackage{helvet}
\renewcommand{\familydefault}{\sfdefault}

\usepackage[margin=18mm]{geometry}
\usepackage{xcolor}
\usepackage{tikz}
\usepackage{tabularx}
\usepackage{enumitem}
\usepackage{tcolorbox}
\usepackage{graphicx}
\usepackage{amssymb}

\definecolor{physik}{HTML}{2c5aa0}
\definecolor{hinweis}{HTML}{b35c00}

\tcbuselibrary{skins,breakable}

\newtcolorbox{merkbox}[1][]{
    colback=white,
    colframe=physik,
    fonttitle=\bfseries,
    title=#1,
    boxrule=1.5pt
}

\newtcolorbox{aufgabebox}[1][]{
    colback=white,
    colframe=hinweis,
    fonttitle=\bfseries,
    title=#1,
    boxrule=1pt
}

% Lücken-Befehl
\newcommand{\luecke}[1]{\underline{\hspace{#1}}}

\pagestyle{empty}

\begin{document}

% ═══════════════════════════════════════════════════════════════════════════
% KOPFZEILE
% ═══════════════════════════════════════════════════════════════════════════

\begin{minipage}{0.6\textwidth}
{\Large\bfseries\color{physik}Elektromagnetische Induktion}\\[2mm]
{\small Physik Klasse 9 | Lernblatt 09b}
\end{minipage}
\hfill
\begin{minipage}{0.35\textwidth}
\raggedleft
Name: \luecke{4cm}\\[2mm]
Datum: \luecke{2.5cm}
\end{minipage}

\vspace{3mm}
\hrule
\vspace{5mm}

% ═══════════════════════════════════════════════════════════════════════════
% KASTEN 1: WAS IST INDUKTION?
% ═══════════════════════════════════════════════════════════════════════════

\begin{merkbox}[Kasten 1: Was ist Induktion?]

\textbf{Definition:}

Bewegt man einen \luecke{2.5cm} relativ zu einer \luecke{2.5cm},

entsteht eine elektrische \luecke{2.5cm}.

Diesen Vorgang nennt man \textbf{elektromagnetische Induktion}.

\vspace{3mm}
\textbf{Wichtig:} Induktion erfordert eine \luecke{3cm}!

Ohne Bewegung (Änderung) gibt es \luecke{2.5cm} Spannung.

\vspace{3mm}
\textbf{Entdecker:} Michael \luecke{2.5cm} (1831)

\end{merkbox}

\vspace{4mm}

% ═══════════════════════════════════════════════════════════════════════════
% KASTEN 2: EINFLUSSFAKTOREN
% ═══════════════════════════════════════════════════════════════════════════

\begin{merkbox}[Kasten 2: Induktionsgesetz (qualitativ)]

\textbf{Die Induktionsspannung ist größer, wenn:}

\begin{enumerate}[noitemsep]
    \item Die Bewegung \luecke{3cm} ist.
    \item Der Magnet \luecke{3cm} ist.
    \item Die Spule mehr \luecke{3cm} hat.
\end{enumerate}

\vspace{3mm}
\textbf{Richtung der Spannung:}

Ändert man die \luecke{3.5cm} der Bewegung,

ändert sich auch das \luecke{3cm} der Spannung (+ $\leftrightarrow$ $-$).

\end{merkbox}

\vspace{4mm}

% ═══════════════════════════════════════════════════════════════════════════
% KASTEN 3: MOTOR UND GENERATOR
% ═══════════════════════════════════════════════════════════════════════════

\begin{merkbox}[Kasten 3: Motor und Generator im Vergleich]

Motor und Generator sind ``\luecke{3cm}'': \textbf{Gleicher Aufbau -- umgekehrte Funktion!}

\vspace{3mm}
\begin{tabularx}{\textwidth}{|X|X|}
\hline
\textbf{Motor} & \textbf{Generator} \\
\hline
Strom fließt \luecke{2cm} & Strom kommt \luecke{2cm} \\[0.6em]
\hline
\luecke{2.5cm} Energie & \luecke{2.5cm} Energie \\
$\downarrow$ & $\downarrow$ \\
Bewegungsenergie & Elektrische Energie \\[0.6em]
\hline
Anwendung: E-Auto, Mixer & Anwendung: \luecke{3cm} \\[0.6em]
\hline
\end{tabularx}

\end{merkbox}

\vspace{5mm}

% ═══════════════════════════════════════════════════════════════════════════
% ÜBUNGEN
% ═══════════════════════════════════════════════════════════════════════════

\begin{aufgabebox}[Übung 1: Anwendungen (4 BE)]

Ordne die Anwendungen dem richtigen Prinzip zu. Kreuze an.

\vspace{2mm}
\begin{tabularx}{\textwidth}{|X|c|c|}
\hline
\textbf{Anwendung} & \textbf{Motor} & \textbf{Generator} \\
\hline
Fahrraddynamo & $\square$ & $\square$ \\[0.4em]
\hline
Elektroauto (Antrieb) & $\square$ & $\square$ \\[0.4em]
\hline
Windkraftanlage & $\square$ & $\square$ \\[0.4em]
\hline
Faraday-Lampe (Schütteltaschenlampe) & $\square$ & $\square$ \\[0.4em]
\hline
\end{tabularx}

\end{aufgabebox}

\vspace{4mm}

\begin{aufgabebox}[Übung 2: Energiekette (3 BE)]

Vervollständige die Energiekette für ein Wasserkraftwerk:

\vspace{3mm}
\begin{center}
\fbox{Lageenergie (Wasser)} $\rightarrow$ \fbox{\luecke{3cm}} $\rightarrow$ \fbox{\luecke{3.5cm}}
\end{center}

\end{aufgabebox}

\vspace{4mm}

\begin{aufgabebox}[Übung 3: Transfer -- 3 BE]

Ein Schüler bewegt einen Magneten langsam durch eine Spule. Das Voltmeter zeigt 0,5~V.

\vspace{2mm}
\textbf{a)} Was passiert mit der Spannung, wenn er den Magneten \textbf{doppelt so schnell} bewegt?

\vspace{1cm}

\textbf{b)} Was passiert, wenn er den Magneten \textbf{anhält}?

\vspace{1cm}

\end{aufgabebox}

\vspace{5mm}

% ═══════════════════════════════════════════════════════════════════════════
% SELBSTCHECK
% ═══════════════════════════════════════════════════════════════════════════

\hrule
\vspace{3mm}

\textbf{Selbstcheck:} Ich kann...

\begin{tabularx}{\textwidth}{|X|c|c|c|}
\hline
& \textbf{sicher} & \textbf{unsicher} & \textbf{noch nicht} \\
\hline
...erklären, was Induktion ist. & $\square$ & $\square$ & $\square$ \\
\hline
...angeben, wovon die Induktionsspannung abhängt. & $\square$ & $\square$ & $\square$ \\
\hline
...Motor und Generator vergleichen. & $\square$ & $\square$ & $\square$ \\
\hline
...Anwendungen der Induktion nennen. & $\square$ & $\square$ & $\square$ \\
\hline
\end{tabularx}

\end{document}
